\clearpage\thispagestyle{empty}
\pdfbookmark[0]{Quellen für Gauß’ Briefe über De Vicos Kometen (Oktober und November 1844)}{cometdevico-sources}
\begin{center}\textbf{Quellen für Gauß’ Briefe über De Vicos Kometen (Oktober und November 1844)}\end{center}
\noindent Diese KI-Übertragung ins heutige Deutsch folgt den Auszügen vom 15. Oktober und 11. November 1844 auf S.~469--470 der \href{https://gdz.sub.uni-goettingen.de/id/PPN236007467}{\textit{Werke}, Band~VI} (Göttingen, 1874). Carl Friedrich Gauß schrieb an Heinrich Christian Schumacher. In den \textit{Astronomischen Nachrichten}, Band~XXII, Nr.~521, \href{https://api.digitale-sammlungen.de/iiif/image/v2/bsb10049447_00147/full/1800,/0/default.jpg}{Sp.~277--278}, steht unter den Auszügen der Name \textit{Gauss}. Die Überschrift der Ephemeride nennt De Vicos Kometen. Das Heft ist laut der Schlusszeile auf \href{https://api.digitale-sammlungen.de/iiif/image/v2/bsb10049447_00148/full/1800,/0/default.jpg}{Sp.~280} auf den 28. November 1844 datiert.
\par\smallskip\noindent Die elliptischen Elemente stammen von Goldschmidt. Der Text in der Werkausgabe endet mit dem Satz, der die Ephemeride ankündigt. Die Zeitschrift druckt sie auf Sp.~277--280 mit einer abschließenden Erläuterung ihrer Spalte zur Lichtstärke. KI-Übertragung und KI-Satz.
\clearpage
